% !TeX root = ../main.tex

% 中文摘要
\begin{abstract}
数字中国战略背景下，教育数字化转型正加速推进学科教学与信息技术的深度融合。针对高中数学实验教学中动态呈现不足、抽象概念可视化程度低及传统教具互动性弱等痛点，本研究创新性地融合动态数学软件GeoGebra与交互式教学平台希沃白板，构建低成本、高交互的数字化实验教学体系。基于建构主义理论、多元智能理论等理论基础，本研究开发出“平面向量基本定理”、“三角函数图象平移伸缩变换”和“圆锥曲线的形成”三个高中数学实验教学资源，涉及课程新授讲解、课堂实操演示和课外兴趣拓展三大模块。依托2019年人教A版教材，深度整合GeoGebra的数学建模优势与希沃白板的课堂交互特性，为发展学生数学建模、直观想象等核心素养提供新型数字化路径，为“双减”背景下课堂增效提质提供可复制的技术方案，有力助推数学教育数字化转型进程。
\end{abstract}



% 英文摘要
\begin{abstract*}
Under the Digital China Strategy, the digital transformation of education is accelerating the deep integration of disciplinary teaching and information technology. Addressing the pain points in high school mathematics experimental teaching—such as insufficient dynamic presentation, low visualization of abstract concepts, and weak interactivity of traditional teaching tools—this study innovatively integrates the dynamic mathematics software GeoGebra and the interactive teaching platform Seewo EasiNote to construct a cost-effective, highly interactive digital experimental teaching system.Guided by constructivism theory and multiple intelligences theory, three experimental teaching resources for high school mathematics are developed in the research: ``the Fundamental Theorem of Plane Vectors'', ``the Translation and Scaling Transformations of Trigonometric Functions'', and ``the Formation of Conic Sections'' covering three key modules: new concept instruction, in-class demonstration, and extracurricular exploration. Building upon the 2019 People's Education Press Version A curriculum, the study synergizes GeoGebra's mathematical modeling strengths with Seewo EasiNote's classroom interactivity, providing a novel digital pathway to develop students' core competencies such as mathematical modeling and visual imagination. It offers replicable technical solutions to enhance classroom efficiency under the ``Double Reduction'' policy, strongly advancing the digital transformation of mathematics education.
\end{abstract*}